\scanpage{47}
catégorie \inserted{additive(et même quasi-abélienne)}avec $\otimes$ qui a toutes les vertus,sauf d'être abélienne, et
sauf d'admettre des \uline{Hom} internes. Pour trouver ces derniers, on peut
si on veut paraphraser la construction de la catégorie des motifs, en
\deleted{introduisant un motif} considérant l'objet \inserted{$\handwritten{\Lambda}(-1)$} de $\tilde{\underline{B}}$ conoyau de
$p^{\handwritten{\unsure{*|\alpha}}}$ : $e \handwritten{\longrightarrow} S$, où $p:S^{1} \handwritten{\longrightarrow} e$ est l'application du cercle sur un
point. C'est un objet de poids 1 (le poids est défini par le degré \deleted{d}
des générateurs, changés de signe). On constate aussitôt que le
foncteur $M\handwritten{\mapsto} M(-1)=M\otimes\handwritten{\Lambda}(-1)$ de $\tilde{\underline{B}}$ dans lui-même est pleinement fidèle,
ce qui permet alors, par une construction formelle standard, de complé%
ter $\tilde{\underline{B}}$ en une catégorie où $\handwritten{\Lambda}(-1)$ devient inversible. Il se trouve
alors que dans cette dernière les \uline{Hom}(X,Y) internes existent et
\deleted{satisfont} sont isomorphes à $\check{X}\otimes Y$ ... en fait , on trouve maintenant la
catégorie de\deleted{s} \uline{tous} les modules gradués libres de type fini sur \quad .
Bien sûr, cette catégorie n'est pas encore abélienne pour autant,
puisque \quad n'est pas un corps; \deleted{pour} \inserted{si on veut} définir la catégorie abélienne
engendrée, on s'attend à trouver l\deleted{\unsure{e}}a\deleted{s} catégorie \handwritten{/}des\handwritten{/}modules de \deleted{type} \inserted{\handwritten{présentation}} fini\inserted{\handwritten{e}}
sur $\handwritten{\Lambda}$.( \inserted{\handwritten{qui est abélienne bien que $\Lambda$ ne soit pas noethérien, car $\Lambda$ est}} \deleted{\quad est-il noethérien, j.e. le nombre de générateurs est-il}
\margin{\handwritten{\unsure{Non|?}}}\deleted{fini ? )\unsure{.}} \handwritten{$\varinjlim$ d'anneaux noeth\unsure{.|,} à hom\unsure{s} de transition \uline{plats}).}

Cette construction ressemblerait à celle d'une théorie des
motifs où on n'aurait pas passé à l'équivalence homologique pour
les cycles, mais où on aurait gardé l'équivalence rationelle, - qui
est bien trop fine \deleted{pou\unsure{r}}à beaucoup d'égards. Pour trouver un équiva\deleted{l}%
lent plus étroit, il convi\deleted{n}e\deleted{t}nt de reprendre la théorie pré\unsure{\deleted{d}c|d}édente,on
pourrait formuler le même Pb universel que plus haut, mais en prenant
l'équivalence homologique (via le graphe) des morphismes de variétés,
au lieu de l'équivalence homotopique. \deleted{Mais alors il semble que l'on
trouve comme catégorie solution celle dont les objets sont les varié-}

\margin{\handwritten{46}}
